\section{Conclusion}

\plat{} is best understood not as a technically novel programming language,
but as a design case study in culturally localized programming-language
design. Its conventional interpreter core is deliberate: by avoiding claims of
new semantic mechanisms, the project brings attention to the cultural and
linguistic layer of language design. Limburgian keywords, diagnostics,
examples, naming choices, documentation, dialect decisions, and ASCII tradeoffs
become the primary subjects of analysis.

The broader lesson is that regional-language programming systems can matter
even when their computational mechanisms are familiar. They can create sites
of visibility, experimentation, teaching, and preservation. Such systems
should be evaluated modestly and carefully, with community feedback and clear
limits, but they deserve a place in conversations about what programming
languages are for and whom they are allowed to represent.
